\scanpage{66}
% block: gerbe class
\[
\handwritten{\xi\in}H^{\handwritten{2}}(k_{\mathrm{pl}},G)
\]
est la classe de la gerbe tannakienne qui définit $\handwritten{\mathcal{M}}$. Donc la connaissance du caractère numérique de $\handwritten{\mathcal{M}}$\deleted{\inserted{(qui implique déjà la donnée}} équivaut à la donnée, en plus \unsure{du|de|da} groupe de t.m. $G$ sur $k$, des classes $u_{\handwritten{\unsure{*|\sigma}}}(\handwritten{\xi})$ pour tous les homomorphismes $u$ de $G$ dans des tores de la forme $\handwritten{\prod}_{Z/k}\underline{G}_m$, où $Z$ est une extension finie de $k$. \deleted{(}On vient de voir en effet que la donnée du caractère numérique permet de reconstituer ce qui précède. Pour voir l'inverse, notons que \deleted{si} le rang $\unsure{\deleted{d}n|\deleted{n}d|n|d}(\handwritten{\unsure{\sigma|\xi}})$ de $M_{\handwritten{\sigma}}$ est déterminé par la formule
\[
n(\handwritten{\sigma})=\operatorname{card}(\handwritten{\sigma})\operatorname{ord}(\handwritten{\xi_\sigma})
\]
(où $\operatorname{card}\handwritten{\sigma}$ est le cardinal de l'orbite $\handwritten{\sigma}$, $\operatorname{ord}(\handwritten{\xi_\sigma})$ est l'ordre de $\handwritten{\xi_\sigma}$ dans $\operatorname{Br}(Z_{\handwritten{\sigma}})$).

% block: determination question
La question de savoir si le caractère numérique détermine à équivalence de catégories tannakiennes près revient donc, pour $G$ fixé, à la question de savoir si un élément \deleted{de} $\handwritten{\xi\in}H^{\handwritten{2}}(k_{\mathrm{pl}},G_{\deleted{\times}})$ est connu quand on connait les $u_{\handwritten{*}}(\handwritten{\xi})$ pour $u:G\handwritten{\unsure{\mathrel{-\! -}|\longrightarrow}}\underline{Z}^{\unsure{\times|x|\mathbb{x}}}$ comme dessus. \deleted{\unsure{On va exes|On va ex}} C'est le cas pour $G$ diagonalisable \deleted{pro-dénombrable} (pro-dénombrable ?) On va examiner deux autres cas.

% block: example 1
\uline{Exemple} \unsure{$\mathbb{I}$|I|1}. $G$ est un tore de dimension 1, donc (s'il n'est trivial) tordu par une extension quadratique $Z$ de $k$. On a donc une suite exacte
\diagram{\begin{tikzcd}[ampersand replacement=\&]
0 \handwritten{\arrow[r]} \& G \handwritten{\arrow[r,dash]} \arrow[r,phantom,"u"] \& \underline{Z}^{\times} \handwritten{\arrow[r]} \arrow[r,phantom,"N_{Z/k}"] \& \underline{G}_m \handwritten{\arrow[r]} \& 0\ ,
\end{tikzcd}}
d'où on conclut, par la suite exacte de cohomologie et le th. 90 que $\operatorname{Ker}(u:H^{\handwritten{2}}(k,G)\handwritten{\longrightarrow}H^{\handwritten{2}}(k,\underline{Z}^{\times}))$ est nul. Donc dans le cas envisagé, le caractère numérique détermine $\handwritten{\mathcal{M}}$ à équivalence près.

% block: cancelled example 2
\deleted{\uline{Exemple} 2. Gardons \deleted{le} $G$ comme dessus, soit $n$ un entier $\quad2$, et soit $G'=\operatorname{Ker}n.\operatorname{id}_G$. La suite exacte}
\diagram{\begin{tikzcd}[ampersand replacement=\&]
\deleted{0} \& \deleted{G'} \& \deleted{G} \arrow[r,phantom,"\deleted{n}"] \& \deleted{G} \& \deleted{0}
\end{tikzcd}}
\deleted{et le résultat précédent montre que l'ensemble des classes de}
% archive sheet number: 65
